\noindent\textbf{Elección del procedimiento.} Para $u=Uy/h$, el gradiente es $U/h$ y la disipación volumétrica $\Phi=\mu(U/h)^2$. Integrar en el espesor da potencia por área $\mu U^2/h$; debe coincidir con el trabajo de la pared por área $\tau U$.

\begin{enumerate}

\item Calcular el esfuerzo constante de Couette.

\[\fitmath{\tau=\frac{\mu\cdot U}{h}}\]

\[\fitmath{\tau=\frac{(0.05)\cdot (3)}{(0.0005)}\approx 300\,\mathrm{Pa}}\]

\item Integrar la disipación volumétrica en el espesor.

\[\fitmath{\mathrm{Parea}=\frac{\mu\cdot {U}^{2}}{h}}\]

\[\fitmath{\mathrm{Parea}=\frac{(0.05)\cdot {(3)}^{2}}{(0.0005)}\approx 900\,\mathrm{W/m^{2}}}\]

\item Contrastar por potencia aportada en la pared móvil.

\[\fitmath{\mathrm{Pwall}=\tau\cdot U}\]

\[\fitmath{\mathrm{Pwall}=(300)\cdot (3)\approx 900\,\mathrm{W/m^{2}}}\]

\end{enumerate}
